% =============================================================================================
\section{Whole-field model, live solver and missing effects}
\label{sec:field}

Sections~\ref{sec:geometry}--\ref{sec:robust_spec} treat one shot straight down the HIVE tipping
plane. This section describes how the same flight model is used for a release point anywhere on the
field (the explorer's Robot view, \code{viz/explorer/solver.js}, with a Python copy in
\code{src/}), which effects were added on top of it, and which effects are still missing. Every
window here is model output\predicted{} with the literature aerodynamics\assumed; none of it is a
hit rate.

% ---------------------------------------------------------------------------------------------
\subsection{Whole-field oblique model}
\label{sec:field-oblique}

The release point is $(x,y)$ in the field frame of the HIVE section: $x$ along the shot axis with the
pivot plane at $x=0$, the far (unfavourable) side at $x<0$, and the up-cell gate centre at
$(x_g,0)$ with $x_g=0.401$\,\unit{m}\derived. The shot is aimed at the gate centre. With
$d=|x|$ the distance to the pivot plane,
\begin{equation}
  L_h=\begin{cases} d+x_g & \text{far side},\\ \max(d-x_g,\,10^{-3}\,\text{m}) & \text{near side},\end{cases}
  \qquad \psi=\arctan\frac{|y|}{L_h},
  \label{eq:field-psi}
\end{equation}
where $\psi$ is the bearing off the HIVE centre line\derived. The ball flies in the vertical plane of
that bearing with its spin axis perpendicular to it, so there is no sideways Magnus force. Its
in-plane path is the planar path of Section~\ref{sec:flight} with the horizontal coordinate and
velocity projected on the shot axis,
\begin{equation}
  X_{\text{axis}}(t)=X(t)\cos\psi,\qquad V_{X,\text{axis}}(t)=V_X(t)\cos\psi,
  \label{eq:field-proj}
\end{equation}
and the lateral offset of the ball centre from the cell centre line is
$s(t)=(x_g-x_{\text{axis}}(t))\tan\psi$. The projection is exact for walls extruded along $y$ and
still air. Bearings above $60^\circ$\assumed{} are flagged as extrapolated.

% ---------------------------------------------------------------------------------------------
\subsection{The 3D cell: side walls and pentagon roof}
\label{sec:field-cell3d}

The precomputed maps use the centre section at one assumed lateral offset per bearing and do not
model the side walls. The live solver instead tests the ball-centre path $(x,y,z)$ against the cell
as 3D faces inflated by $R=0.0355$\,\unit{m}\confirmed: the centre section extruded across the cell,
side walls at $\pm0.254$\,\unit{m}\confirmed{} with inner face, rim and outer face, wall thickness
$0.0508$\,\unit{m}\assumed, and a pentagon roof whose ridge drops linearly towards the side walls
(from the opening profile $0.3556\to0.19329$\,\unit{m} over $0.254$\,\unit{m}\derived). The
clearance is the exact segment-to-polygon distance minus $R$. A shot is clean when the centre enters
the cell interior before any face comes closer than $R$.

At $\psi=0$ with the ridge off the 3D test reproduces the section test exactly (2500 shots per side).
With the ridge on it is slightly stricter on the centre line, because the roof is lower off it. Off
the centre line the 3D cell is usually more permissive than the old single-offset section: at
$d=1.0$\,\unit{m}, $h=0.5$\,\unit{m}, $S=1.0$ the far-side clean grid cells go from $306$ to $301$ at
$\psi=0$, from $212$ to $309$ at $20^\circ$ and from $164$ to $362$ at $40^\circ$\predicted{} (no-spin
twin $348\to346$, $253\to347$, $186\to401$\predicted). These are counts on a coarse search grid, not
probabilities.

% ---------------------------------------------------------------------------------------------
\subsection{Live solver: knife-edge and tolerant minima}
\label{sec:field-solver}

For a continuous release point and height the solver flies a coarse grid (exit speed $2$--$16$\,\unit{m/s}
in $0.1$\,\unit{m/s} steps, hood angle in $0.1^\circ$ far-side or $0.25^\circ$ near-side steps) and then
refines near the minima. It reports two kinds of minimum:
\begin{itemize}
  \item \textbf{Knife edge} (\code{minRpm}, \code{minAngle}): the lowest clean exit speed or hood angle the
    search found. These are grazing shots. Their clearance is typically $0$--$0.5$\,\unit{mm}\predicted, and the
    value moves with the resolution because the window edge carries about $0.1^\circ$ of path-sampling noise
    at the $2$\,\unit{ms} step. They are shown, but a robot must not be set there.
  \item \textbf{Tolerant} (\code{minRpmTol}, \code{minAngleTol}): the lowest speed whose continuous clean
    hood-angle run is at least $0.2^\circ$\assumed{} wide, and the lowest hood angle whose continuous clean
    speed run is at least $0.1$\,\unit{m/s}\assumed{} wide, with clearance $\ge2$\,\unit{mm}\assumed{} at every
    sample. The most forgiving shot is the centre of the largest all-clean square on the search grid.
\end{itemize}
Every result is solved twice: with the chosen backspin and as the no-spin twin ($S=0$, no Magnus
term), with the same air and drag settings. The fused integrate-and-classify fast path is checked
against the reference integrator and classifier for equal class and clearance to $10^{-12}$\,\unit{m}.

At the explorer's default robot spot ($x=-1.0$, $y=0.2$\,\unit{m}\userval, $h=0.737$\,\unit{m}\userval,
$\psi=8.1^\circ$\derived, $S=1.0$\assumed, $v/u=0.42$\assumed, nominal aerodynamics) the solver gives
\begin{center}
\begin{tabular}{@{}lcc@{}}
  \toprule
  & backspin $S=1.0$ & no-spin twin $S=0$ \\
  \midrule
  knife-edge min.\ rpm & $2677$\,rpm at $71.74^\circ$\predicted{} & $2682$\,rpm at $72.86^\circ$\predicted \\
  tolerant min.\ rpm & $2784$\,rpm at $73.30^\circ$\predicted{} & $2784$\,rpm at $74.47^\circ$\predicted \\
  tolerant min.\ hood angle & $75.80^\circ$ at $3023$\,rpm\predicted{} & $77.30^\circ$ at $3034$\,rpm\predicted \\
  most forgiving & $4269$\,rpm at $80.8^\circ$\predicted{} & $4872$\,rpm at $84.7^\circ$\predicted \\
  \bottomrule
\end{tabular}
\end{center}
The rpm values are derived from the exit speed with the assumed grip and a $98$\,\unit{mm}\userval{}
wheel. The knife-edge clearance at the minimum rpm is $0.63$\,\unit{mm}\predicted. At this spot backspin
does not lower the tolerant minimum wheel speed; it lowers the hood angle at that speed by about
$1.2^\circ$\predicted. More spin is not always better: the lift saturates (Section~\ref{sec:flight}).

% ---------------------------------------------------------------------------------------------
\subsection{Robot and turret view}
\label{sec:field-robot}

The Robot view places a $0.457$\,\unit{m}\confirmed{} robot anywhere on the field at any heading. The
release point is the robot centre at the chosen height\assumed, and a turret aims at the cell centre
exactly\assumed. The page reports the turret yaw relative to the robot front and, if the chosen turret
travel ($360^\circ$, $\pm180^\circ$, $\pm90^\circ$, $\pm45^\circ$ or fixed) cannot reach that yaw, how far
to turn the robot. Poses inside the A-frame base are illegal; the second HIVE is only flagged as an
obstacle. The recommended setting is the most forgiving shot, unless it needs more exit speed than the
grip limit below. In that case the lowest tolerant wheel speed is recommended. At the default spot the most forgiving
shot ($4269$\,rpm, about $9.2$\,\unit{m/s}\predicted) is above the limit, so the page recommends
$2784$\,rpm at $73.3^\circ$\predicted{} (twin $2784$\,rpm at $74.5^\circ$\predicted).

% ---------------------------------------------------------------------------------------------
\subsection{Venue air, drag options, drafts and the grip limit}
\label{sec:field-options}

These were added from the ranked list in \code{research/MISSING\_EFFECTS.md} because they need no new
measurement. Python (\code{src/air.py}, \code{src/feasibility.py}) and JavaScript (\code{physics.js},
\code{solver.js}) carry the same formulas. Tests compare trajectories from the two at
$2\times10^{-6}$ relative, which is the float32 storage limit.

\paragraph{Air density.} Simplified CIPM-2007 with a Magnus-form vapour pressure and a US Standard
Atmosphere pressure from altitude,
\begin{equation}
  \rho=\frac{p-0.378\,e}{R_d\,T},\quad e=\phi\,611.2\exp\frac{17.62\,t}{243.12+t},\quad
  p=101325\,(1-2.25577\times10^{-5}\,z)^{5.25588},
  \label{eq:field-rho}
\end{equation}
with $R_d=287.06$\,\unit{J/(kg\,K)}, $t$ in \unit{$^\circ$C}, $T=t+273.15$, humidity $\phi$ and altitude
$z$ in \unit{m} (user inputs). Viscosity follows Sutherland's law. Still air at $20^\circ$C,
$101.325$\,\unit{kPa}, dry gives $\rho=1.2041$\,\unit{kg/m\textsuperscript{3}}\derived{} and
$\nu=1.506\times10^{-5}$\,\unit{m$^2$/s}\derived. This is now the live solver's default instead of the
rounded $1.20$\assumed{} used by the precomputed maps. At the default spot the change moves the tolerant
minimum from $2793$\,rpm at $73.43^\circ$ to $2784$\,rpm at $73.30^\circ$\predicted. The twin stays at
$2784$\,rpm at $74.47^\circ$\predicted{} to the precision shown.

\paragraph{Drag options.} $C_d$ stays the constant literature value by default. Two alternatives are
available. The first is the smooth-sphere correlation $C_d(\mathrm{Re})$ of Morrison (2013), a reference
curve for a ball without holes and not a POLLEN fact. The second is a what-if partial drag crisis that
falls from the constant $C_d$ to $0.33$\assumed{} with a smoothstep between $\mathrm{Re}=3\times10^4$ and
$5\times10^4$\assumed. Backspin can add drag:
\begin{equation}
  C_d\leftarrow C_d+\Delta C_d\,\min(S/0.6,\,1),\qquad \Delta C_d\in\{0,\,0.08,\,0.18\}\assumed .
  \label{eq:field-spindrag}
\end{equation}
The no-spin twin has $S=0$ and never gets this term.

\paragraph{Draft.} A steady draft $\vect{w}=(w_\parallel,w_z)$ in the shot plane (user input) enters
through the relative velocity $\vect{v}_{\text{rel}}=\vect{v}-\vect{w}$. Drag, Magnus force and $S$ all
use it. A sideways part is an option on the 3D path (Section~\ref{sec:field-move}).

\paragraph{Grip limit.} An answer that needs an exit speed above $8.9$\,\unit{m/s}\predicted{}
(the hood launcher model of Section~\ref{sec:launcher}) is flagged, not removed. Above that speed,
$v/u$ and $S$ are uncertain.

\paragraph{Effect at the default spot} (tolerant minimum rpm and its hood angle, twin in brackets; all
\predicted):
\begin{center}
\begin{tabular}{@{}lll@{}}
  \toprule
  option & backspin $S=1.0$ & twin $S=0$ \\
  \midrule
  still air (default) & $2784$\,rpm, $73.30^\circ$ & $2784$\,rpm, $74.47^\circ$ \\
  Morrison $C_d(\mathrm{Re})$ & $2775$\,rpm, $73.64^\circ$ & $2765$\,rpm \\
  what-if drag drop & as still air ($\mathrm{Re}<3\times10^4$ here) & as still air \\
  spin drag $\Delta C_d=0.18$ & $2821$\,rpm, $72.98^\circ$ & $2784$\,rpm (no term) \\
  tail draft $0.3$\,\unit{m/s} & $2803$\,rpm, $73.84^\circ$ & $2784$\,rpm, $74.8^\circ$ \\
  updraft $0.3$\,\unit{m/s} & $2793$\,rpm, $73.84^\circ$ & $2765$\,rpm \\
  altitude $1600$\,\unit{m} ($\rho=0.9925$\derived) & $2775$\,rpm, $73.77^\circ$ & $2765$\,rpm, $74.61^\circ$ \\
  \bottomrule
\end{tabular}
\end{center}
No tolerant minimum at this spot is grip-limited. Each option moves the hood angle by up to about
$0.5^\circ$\predicted, which is wider than the $0.2^\circ$\assumed{} tolerance. A venue therefore needs
its own aim check.

% ---------------------------------------------------------------------------------------------
\subsection{Shooting on the move, hit chance and lift options}
\label{sec:field-move}

These options come from the same ranked list (\code{research/MISSING\_EFFECTS.md}, build-next 2a--2f). All
are off by default, so the default answer above is unchanged. Python (\code{src/robustness.py},
\code{src/air.py}, \code{src/flight.py}) and JavaScript (\code{solver.js}, \code{physics.js}) agree on
launch states to $10^{-12}$, on 3D paths to $10^{-9}$ relative and on every Monte Carlo count
(\code{tests/test\_move.py}, \code{tests/js\_move\_check.js}).

\paragraph{Moving robot.} The robot's ground velocity at the exit point is
$\vect{v}_r=a\,\hat{\vect{u}}+c\,\hat{\vect{n}}+\omega_z\hat{\vect{z}}\times\vect{r}_{\text{exit}}$,
with $\hat{\vect{u}}$ the horizontal aim direction, $\hat{\vect{n}}=\hat{\vect{z}}\times\hat{\vect{u}}$, and the
speeds $a$, $c$ and turn rate $\omega_z$ user inputs (field-frame $v_x$, $v_y$ are also accepted). The ball
leaves with
\begin{equation}
  \vect{v}_0=\vect{v}_{\text{shooter}}(v,\ \theta_e)+\vect{v}_r,\qquad \theta_e=\theta_{\text{hood}}+\delta_0,
  \label{eq:field-move}
\end{equation}
where $\delta_0$ is an exit-angle offset (\assumed, default $0$). With lead on (default) the shooter turns so the
sideways part of $\vect{v}_0$ cancels, which needs $|\vect{v}_r\cdot\hat{\vect{n}}|\le v\cos\theta_e$. The
backspin axis is horizontal and normal to the shooter direction, plus $\omega_z$ about the vertical; the
no-spin twin has no backspin. The flight is then 3D in the 3D cell, with a sideways draft allowed. The solver
also reports how far the centre of the clean run moves per $0.1$\,m/s toward the goal and sideways.

\paragraph{Hit chance.} At one setting the solver fires $n=1000$ shots per aero set (lift low, nominal and
high) with ASSUMED Gaussian launch scatter: speed $1.5\,\%$, elevation $0.40^\circ$, heading $0.50^\circ$
($1\sigma$, the nominal set of Section~\ref{sec:robust_spec-mc}). All sets and the twin use the same seeded
draws. The result is $k/n$ with a Wilson $95\,\%$ interval (Section~\ref{sec:robust_spec-wilson}). It is a
model probability under assumed scatter, not a measured repeatability. At the default spot and the recommended
$2784$\,rpm, $73.3^\circ$: nominal $0.124$ $[0.105,\,0.146]$, lift low $0.092$, lift high $0.074$; the twin
at its own setting ($2784$\,rpm, $74.47^\circ$) about $0.13$ (all \predicted). Backspin does not raise the
chance here.

\paragraph{Lift options.} Gyro-correct spin ratio $S_\perp=|\vect{\omega}\times\vect{v}_{\text{rel}}|R/
|\vect{v}_{\text{rel}}|^2$ on the 3D path (equal to $S$ for pure backspin, tested). A user table $C_l(S)$,
linear between points and flat after the last. Reverse Magnus: when $\mathrm{Re}<4.6\times10^4$ and
$S<0.5$ the lift blends (smoothstep near both bounds) to $-C_{l,\text{rev}}\min(S/0.1,1)$, $C_{l,\text{rev}}\in[0,0.1]$\assumed{} (smooth-sphere
data of Kim et al.\ 2014 and Kray et al.\ 2012, not POLLEN).

\paragraph{Effective mass.} $\varepsilon=(0.5+\phi)\rho V/m$ with $\phi\in[0,1]$\assumed{} the share of the
air inside the ball that moves with it, and
$\dot{\vect{v}}=(\vect{F}_{\text{aero}}/m+\vect{g})/(1+\varepsilon)$. For POLLEN
$\varepsilon=0.0045$ to $0.0136$\derived. Default $\varepsilon=0$.

\paragraph{Effect at the default spot} (tolerant minimum rpm and hood angle; \predicted):
\begin{center}
\begin{tabular}{@{}lll@{}}
  \toprule
  option & backspin $S=1.0$ & no-spin twin \\
  \midrule
  none (default) & $2784$\,rpm, $73.30^\circ$ & $2784$\,rpm, $74.47^\circ$ \\
  moving, all speeds $0$ & same & same \\
  $0.5$\,m/s toward the goal & $2738$\,rpm, $78.09^\circ$ & $2719$\,rpm, $79.01^\circ$ \\
  $0.3$\,m/s sideways, lead & $2812$\,rpm, $73.41^\circ$ & $2784$\,rpm, $74.15^\circ$ \\
  $0.3$\,m/s sideways, no lead & no clean shot & no clean shot \\
  turning $90^\circ$/s & $2793$\,rpm, $73.42^\circ$ & as default \\
  $\delta_0=+1^\circ$ & $2784$\,rpm, $72.30^\circ$ & $2784$\,rpm, $73.47^\circ$ \\
  sideways draft $0.3$\,m/s & $2803$\,rpm, $73.54^\circ$ & $2793$\,rpm, $74.61^\circ$ \\
  $C_l$ table $(0.5,0.2),(1.5,0.35)$ & $2803$\,rpm, $73.25^\circ$ & as default \\
  $\phi=1$ & $2765$\,rpm, $73.34^\circ$ & $2756$\,rpm, $74.36^\circ$ \\
  \bottomrule
\end{tabular}
\end{center}
Driving toward the goal moves the hood angle by about $1^\circ$ per $0.1$\,m/s\predicted, several times the
ASSUMED aim tolerance. Sideways motion with lead halves the backspin window (315 to 158 clean pairs, twin 364
to 347), because the spin axis follows the turned shooter. Reverse Magnus does not act here ($S=1$).

% ---------------------------------------------------------------------------------------------
\subsection{Hit-chance guidance and the measurement protocol}
\label{sec:field-chance}

\paragraph{Best-chance setting.} The lowest-rpm setting above sits near the edge of the clean window, so
the page now recommends the clean setting with the highest model chance to go in instead
(\code{solver.js} \code{bestChance}, replayed in \code{src/robustness.py} \code{best\_chance}). The
candidates are clean cells inside the PREDICTED grip limit ($8.9$\,m/s), ranked by their share of clean
neighbours under a Gaussian with the scatter as $\sigma$ and picked at least one $\sigma$ apart (up to $24$),
plus the tolerant minima. Stage~1 scores each with $64$ shots per air set; at most $5$ whose pooled Wilson
upper bound reaches the best lower bound go to stage~2 with $400$ shots. All candidates in a stage see the
same seeded draws (common random numbers). The winner is scored again on $1000$ fresh shots per air set, so
the reported chance is not biased by the selection. The score is the mean over the three air sets (lift low,
nominal, high), equal weight\assumed. The no-spin twin is searched on its own grid, and one scatter source at a
time names the one that matters most.

\paragraph{Default spot} ($x=-1.0$, $y=0.2$\,m, $h=0.737$\,m\userval, $S=1.0$\userval, grip $0.42$\assumed,
scatter $1.5\,\%$, $0.40^\circ$, $0.50^\circ$\assumed; all chances \predicted, Wilson $95\,\%$ ranges):
\begin{center}
\begin{tabular}{@{}lllll@{}}
  \toprule
  setting & lift low & nominal & lift high & mean \\
  \midrule
  best chance, $3062$\,rpm\derived, $76.2^\circ$ & $0.207$ & $0.295$ $[0.268,\,0.324]$ & $0.064$ & $0.189$ \\
  lowest rpm, $2784$\,rpm\derived, $73.3^\circ$ & $0.092$ & $0.124$ $[0.105,\,0.146]$ & $0.074$ & $0.097$ \\
  no-spin twin at its best, $3851$\,rpm\derived, $82.1^\circ$ & \multicolumn{3}{l}{$0.375$ $[0.346,\,0.405]$ (no lift: one set)} & $0.375$ \\
  no-spin twin at its lowest rpm, $2784$\,rpm, $74.47^\circ$ & \multicolumn{3}{l}{$0.132$ $[0.112,\,0.154]$} & $0.132$ \\
  \bottomrule
\end{tabular}
\end{center}
Here the hood angle is the scatter source that matters most for both shots: with only one source on, the
nominal chance is $0.46$ (speed only), $0.40$ (angle only) and $1.00$ (aim only)\predicted; for the twin
$0.74$, $0.44$ and $1.00$\predicted. At $3\,\%$ speed scatter\userval{} the best setting moves to $3341$\,rpm,
$78.0^\circ$ (nominal $0.258$), speed matters most, and the twin's best is $3990$\,rpm, $82.6^\circ$ ($0.290$)
(all \predicted). At this spot the no-spin twin has the higher chance: more spin is not always better. Over the
lift range the backspin chance moves by a factor of almost five ($0.064$ to $0.295$), so the measured $C_l$
decides whether backspin helps here at all.

\paragraph{What the chance is not.} It is a model probability under ASSUMED scatter and ASSUMED air sets, not
a measured repeatability. It becomes a prediction a team can test only after the scatter itself is measured
(the standard deviation of exit speed and elevation over about $20$ shots at one setting, logged with wheel
speed and battery voltage) and the hit rate is counted with a Wilson interval from many shots
(Section~\ref{sec:robust_spec-wilson}).

\paragraph{Measurement protocol.} \code{docs/MEASURE.md} and the browser kit \code{viz/measure/measure.html}
(linked from the explorer) turn the two largest unknowns into numbers with one phone camera at
$240$\,fps or faster. The ball centre is clicked frame by frame and the track is fitted with the same RK4
flight model as \code{physics.js} (Python \code{src/fit\_aero.py} agrees with the page to $5\times10^{-9}$).
\begin{itemize}
  \item \textbf{Drop test} (no launcher, no spin): $C_d$ at zero spin. With one $4$\,m drop and $4$\,mm click
    noise the $95\,\%$ error is about $\pm0.008$, and about $\pm0.085$ from $1.5$\,m\derived.
  \item \textbf{Twin test:} shot~A with no spin gives $v_0$, the launch angle and $C_d$; shot~B, the same
    setting with backspin, holds A's $C_d$ and gives $C_l$ beside its no-spin twin. On synthetic
    \code{physics.js} tracks (truth ASSUMED, $C_d=0.50$, $C_l=0.25$) one pair gives sd $0.008$ in $C_d$ and
    $0.010$ in $C_l$ at $4$\,mm noise and $0.7$\,s tracked\derived. These are synthetic checks, not POLLEN data.
  \item \textbf{Length scale.} The scale bar is the largest error: $+1\,\%$ in scale moves $C_l$ by about
    $19\,\%$\derived. Shot~A also fits the scale against gravity (sd about $0.3\,\%$\derived), and the fit can
    use either scale.
  \item \textbf{Exit speed and spin:} wheel rpm with the video gives the grip $v/u$; counting turns of a
    striped ball gives $S$ at the first tracked frame.
\end{itemize}
The results go into the explorer's Advanced settings (nearest air preset and grip), and the best-chance search
then runs on measured rather than assumed inputs. Until then every chance above stays a model prediction.

% ---------------------------------------------------------------------------------------------
\subsection{Missing effects and how to measure them}
\label{sec:field-missing}

\code{research/MISSING\_EFFECTS.md} ranks 32 effects by how far they move the answer. The largest ones
that are still not in the model, or are in it only as bracket sets, need measurements rather than code:
\begin{center}
\begin{tabular}{@{}p{0.36\linewidth}p{0.58\linewidth}@{}}
  \toprule
  effect & how to measure \\
  \midrule
  Spin lift level $C_l(S)$ (slope and cap) & side video at 240--480 fps, wheel above vs below at equal rpm, taped-hole ball (T1) \\
  Drag level $C_d$ & the same flight video and a drop test (T1) \\
  Grip-limited exit speed, slip $v/u$, ball squeeze & exit video with a wheel encoder (T2); mass, diameter, squeeze (T0/T4) \\
  Hood angle $\ne$ exit angle, exit scatter & exit video over a hood-angle sweep (T2) \\
  Knuckle force on a low-spin holed ball & shot scatter with and without spin (T1) \\
  Exit skew and spin-axis tilt (sideways error) & axle alignment check; spin axis from video \\
  Balls and NECTAR in the cell, restitution, TIP mid-volley & drop tests on skin, rib and NECTAR; cell tests (T5/T6) \\
  Effective mass (an ASSUMED option only) & absorbed by tuning to real shots; mass and volume (T0) \\
  Flywheel droop, battery; real scatter while moving & burst tests with encoder and voltage logs; shots on the move \\
  Venue drafts and air state & air speed at robot height, temperature, humidity, station pressure \\
  \bottomrule
\end{tabular}
\end{center}
Until those numbers exist, the model's job is to say which measurement changes the answer most. It
does not predict a scoring rate.
